% !TEX program = xelatex
% Diese Datei MUSS mit XeLaTeX (oder LuaLaTeX) kompiliert werden — nicht mit
% pdfLaTeX. Grund: fontspec + polyglossia (s. u.) verlangen XeTeX/LuaTeX.
\documentclass[10pt]{article}

% --- Page: A5, with a 1.5cm top stub reserved for the perforated ballot-number band ---
\usepackage[paper=a5paper, top=24mm, bottom=22mm, left=11mm, right=11mm]{geometry}

\usepackage{fontspec}
\usepackage{polyglossia}
\setmainlanguage{german}

\usepackage{amssymb}      % math symbols
\usepackage{tcolorbox}    % boxed questions
\usepackage{array}        % ragged-right p{} columns
\usepackage{setspace}
\setstretch{1.0}
\usepackage{eso-pic}      % page-anchored background layer (fiducials, perforation, QR)
\usepackage{tikz}
\usepackage{qrcode}

\pagestyle{empty}

% 5-digit ballot number for this print run; starts with 5, links the slip to the canvasser record (see PAP "Extensions")
\newcommand{\ballotnumber}{50001}
\newcommand{\ballotband}{\ballotnumber\quad--\quad\ballotnumber\quad--\quad\ballotnumber}

\tcbset{colback=white, colframe=black!55, boxrule=0.4pt, sharp corners,
        left=7pt, right=7pt, top=3pt, bottom=3pt, before skip=5pt, after skip=5pt}
\newcommand{\fragebox}[1]{\begin{tcolorbox}#1\end{tcolorbox}}

% Checkbox: a drawn square, a bit larger than the old math \Box, vertically centred on the
% text line (via \raisebox) so the adjacent number sits on the SAME line, not as a subscript.
\newcommand{\boxsize}{3.6mm}
\newcommand{\ckbox}{\raisebox{-0.2em}{\tikz{\draw[line width=0.5pt] (0,0) rectangle (\boxsize,\boxsize);}}}

% A 5-point Likert row: endpoint labels ABOVE the scale, in narrow wrapping columns tied to
% the extreme boxes (so they never sit above the 2nd / 2nd-last number); then the checkbox row.
\newcommand{\likertrow}[2]{%
  \par\vspace{8pt}
  \noindent
  \parbox[b]{0.30\linewidth}{\raggedright\small #1}%
  \hfill
  \parbox[b]{0.30\linewidth}{\raggedleft\small #2}%
  \par\vspace{2pt}
  \begin{tabular*}{\linewidth}{@{\extracolsep{\fill}}ccccc}
    \ckbox~1 & \ckbox~2 & \ckbox~3 & \ckbox~4 & \ckbox~5
  \end{tabular*}%
}

% --- Background layer: corner fiducials, perforation line, ballot-number band, QR code ---
% All four fiducials sit BELOW the perforation line, i.e. in the part of the page that
% remains once the top stub (with the ballot-number band) is torn off.
\AddToShipoutPictureBG{%
  \begin{tikzpicture}[remember picture, overlay]
    % OCR/OMR registration fiducials: 5mm filled squares
    % -- bottom corners: 3mm inset from the physical corner
    \fill[black] ([xshift=3mm,yshift=3mm]current page.south west) rectangle ++(5mm,5mm);
    \fill[black] ([xshift=-8mm,yshift=3mm]current page.south east) rectangle ++(5mm,5mm);
    % -- top corners: moved down so they sit 3mm below the perforation line, not in the stub
    \fill[black] ([xshift=3mm,yshift=-18mm]current page.north west) rectangle ++(5mm,-5mm);
    \fill[black] ([xshift=-8mm,yshift=-18mm]current page.north east) rectangle ++(5mm,-5mm);

    % Perforation: dashed line 15mm below the top edge, full page width
    \draw[dashed, thick] ([yshift=-15mm]current page.north west) -- ([yshift=-15mm]current page.north east);

    % Ballot-number band, repeated, centered in the top stub
    \node at ([yshift=-7.5mm]current page.north) {\large\ballotband};

    % Fold line: marks the MIDDLE of the sheet that remains after the numbered perforation
    % strip is torn off. Remaining sheet spans 15mm..210mm -> midpoint at 112.5mm from top.
    % It is placed so it falls in the gap between question 2 and question 3 (not through a box).
    \draw[dashed, line width=0.5pt] ([yshift=-112.5mm]current page.north west) -- ([yshift=-112.5mm]current page.north east);
    \node[fill=white, inner sep=3pt] at ([yshift=-112.5mm]current page.north) {\small\itshape hier falten};

    % QR code, bottom-left, clear of the fiducial square
    \node[anchor=south west, inner sep=0pt] at ([xshift=12mm,yshift=10mm]current page.south west) {\qrcode[height=1.4cm]{\ballotnumber}};
  \end{tikzpicture}%
}

\begin{document}

\begin{center}
    {\Large\bfseries Was denken Sie\dots?}\\[1.5mm]
    {\normalsize Ihre Antworten sind völlig anonym und lassen sich nicht zu Ihnen zurückverfolgen.}
\end{center}

\vspace{9mm}

\fragebox{
    \textbf{1.} Wie wahrscheinlich ist es, dass Sie Bündnis~90/Die~Grünen bei den Berliner Wahlen 2026 wählen werden?

    \likertrow{Sehr\\unwahrscheinlich}{Sehr\\wahrscheinlich}
}

\fragebox{
    \textbf{2.} Was würden Sie sagen, wie sehr interessieren sich die Grünen für die Meinung von Menschen wie Ihnen?

    \likertrow{Interessieren sich\\überhaupt nicht}{Interessieren sich\\sehr stark}
}

% --- Lücke für die Faltlinie ---
% Die sichtbare Faltlinie wird in der Hintergrundebene exakt bei 112,5mm gezeichnet
% (Mitte des Blatts nach Abtrennen des Nummern-Streifens). Diese Lücke sorgt nur dafür,
% dass zwischen Frage 2 und Frage 3 genug Freiraum bleibt, damit keine Box die Linie kreuzt.
\vspace{20mm}

\fragebox{
    \textbf{3.} Welche Koalition würden Sie sich nach der Wahl in Berlin wünschen?

    \vspace{4pt}
    \setlength{\tabcolsep}{2pt}%
    \begin{tabular}{@{}>{\raggedright\arraybackslash}p{0.43\linewidth} >{\raggedright\arraybackslash}p{0.55\linewidth}@{}}
        \ckbox~,,GroKo`` \mbox{\footnotesize(CDU, SPD)}           & \ckbox~,,Rot-Rot-Grün`` \mbox{\footnotesize(Grüne, Linke und SPD)} \\[4pt]
        \ckbox~,,Kenia`` \mbox{\footnotesize(CDU, Grüne und SPD)} & \ckbox~,,Schwarz-Grün`` \mbox{\footnotesize(CDU und Grüne)}        \\[4pt]
        \ckbox~Andere: \makebox[2.2cm]{\hrulefill}               & \ckbox~Weiß ich nicht                                             \\
    \end{tabular}
}

\fragebox{
    \textbf{4.} Wie hat Ihnen das Haustürgespräch gefallen?

    \likertrow{Sehr schlecht}{Sehr gut}
}

\vspace{3mm}
\begin{center}
    {\small Bitte falten Sie den Fragebogen in der Mitte und werfen Sie ihn in die Urne.}
\end{center}

\end{document}
